\section{Experimental Methodology}
\label{sec:exp}

\subsection{Evidence Source and Data Audit}

All quantitative results in this paper come from the immutable review package
\texttt{results/measured\_20261014}. The run records source commit
\texttt{4bd07c6}, an empty working-tree status, and the input-CSV SHA-256 digest
\texttt{c733cc0c3e0950fa1d05d57d34c69d0a}\-
\texttt{034383e0c1a2fde0e6878b9e2e8c3668}. Hashes establish byte identity, not
factual provenance.

The table contains 2,000 rows, 2,381 feature columns, and one binary label. The
audit found no exact duplicate vectors. A deterministic stratified row split with
seed 42 assigned 60\% to the training/pool partition, 20\% to validation, and
20\% to test. Table~\ref{tab:data} reports the class counts. No original file hash,
collection time, family, source, or label-definition manifest is available. We
therefore do not claim that the positive class exclusively represents ransomware
or that the rows form a verified subset of SOREL-20M or EMBER.

\begin{table}[t]
\centering
\caption{Audited pilot partitions.}
\label{tab:data}
\small
\begin{tabular}{lrrr}
\toprule
Partition & Negative & Positive & Total \\
\midrule
Training/pool & 900 & 300 & 1,200 \\
Validation & 300 & 100 & 400 \\
Test & 300 & 100 & 400 \\
\midrule
Total & 1,500 & 500 & 2,000 \\
\bottomrule
\end{tabular}
\end{table}

The dimensional layout matches the EMBER v2 extractor reference. The groups in
Table~\ref{tab:feature-groups} describe an expected schema, not proof that the
CSV was generated from the documented extractor. Hashed imports, exports, and
section attributes can contain collisions; an individual coordinate need not map
to one human-readable API or semantic property.

\begin{table}[t]
\centering
\caption{Reference feature-group dimensions for a 2,381-coordinate EMBER v2 representation.}
\label{tab:feature-groups}
\small
\begin{tabular}{lr}
\toprule
Feature group & Dimensions \\
\midrule
Byte histogram & 256 \\
Byte-entropy histogram & 256 \\
String extractor & 104 \\
General file information & 10 \\
Header information & 62 \\
Section information & 255 \\
Imports & 1,280 \\
Exports & 128 \\
Data directories & 30 \\
\midrule
Total & 2,381 \\
\bottomrule
\end{tabular}
\end{table}

\subsection{Active-Learning Protocol}

We used five seeds: 11, 22, 33, 44, and 55. For each seed, stratified sampling
selected 200 initially labeled training rows; all five acquisition policies shared
that exact initial set and remaining pool. Each run completed five acquisition
rounds of 100 labels, producing budgets of 200, 300, 400, 500, 600, and 700
training labels. The 400 validation labels are additional and are not counted in
the training budget.

At every round and for every method, a new LightGBM booster was trained with 60
trees, learning rate 0.05, 15 leaves, maximum depth 4, and minimum 10 samples per
leaf. Feature and bagging fractions were one, and deterministic column-wise
training used two threads. Method-specific random, data, and feature seeds were
locked to the run seed. The compared strategies were Random, Uncertainty, $U+H$,
$U+V$, and $\xfull$ with weights specified in Section~\ref{sec:method}.

The threshold was selected from validation predictions to target TPR 0.95. We
then saved validation and test probabilities, predictions, labels, row IDs,
thresholds, and round budgets. For queried rows we saved rank, composite score,
raw $U/H/V$ signals, normalized signals, and revealed label. This enables both
selection replay and checks that no test or validation index enters training.

The primary ranking metrics were ROC--AUC and AP. Thresholded metrics were FPR
and TPR. We report mean and sample standard deviation across seeds, paired
within-seed differences, and normalized area under the learning curve (AULC)
computed by trapezoidal integration over the training-label budget. We do not use
five seed values as five independent test datasets. No target validation AUC was
predeclared, so the run does not estimate labels saved to a target. Query timing
includes pool prediction, TreeSHAP, additivity checking, and normalization; it
excludes labeling, training, sorting, and file I/O.

\begin{table}[t]
\centering
\caption{Core experimental configurations.}
\label{tab:configuration}
\small
\begin{tabular}{p{2.4cm}p{4.6cm}}
\toprule
Component & Configuration \\
\midrule
AL budgets & 200 initial; $5\times100$ queried; 700 final \\
AL seeds & 11, 22, 33, 44, 55 \\
AL model & 60 trees; depth 4; 15 leaves; $\eta=0.05$ \\
AL threshold & Validation-selected at target TPR 0.95 \\
Fidelity model & 100 trees; depth 4; 15 leaves; $\eta=0.1$; seed 42 \\
Fidelity masks & $k\in\{10,25,50\}$; 20 random controls per $k$ \\
Fidelity bootstrap & 500 stratified paired test resamples \\
Execution & Python 3.12.14; LightGBM 4.7.0; two threads \\
\bottomrule
\end{tabular}
\end{table}

\subsection{Fidelity Protocol}

The fidelity experiment trained a separate LightGBM classifier on all 1,200
training rows with 100 estimators, learning rate 0.1, 15 leaves, maximum depth 4,
minimum child size 20, random seed 42, and two threads. It evaluated all 400 test
rows. This model is not an active-learning final-round model; AUC values from the
two experiments must not be subtracted or interpreted as component-wise gains.

We computed native TreeSHAP contributions and checked raw-score additivity. For
each test sample and $k\in\{10,25,50\}$, the $k$ highest absolute contributions
were replaced using the training-column means. Twenty sample-specific uniform
random masks without replacement formed the control at each $k$. We recorded
the selected indices, number of coordinates that actually changed, original and
masked probabilities, AUC, AUC drop, signed and absolute probability changes,
and $\mathrm{FF}_{\mathrm{raw}}$. Random-control summaries report means and
sample SD across masks.

For uncertainty in the top-SHAP AUC drop, 500 stratified bootstrap resamples
preserved the original/masked prediction pair for each sampled test row. These
intervals are conditional on the one fitted model; they do not include
between-training variability. EC and CF were calculated once from the fixed test
explanations according to (\ref{eq:ec}) and (\ref{eq:cf}).

\subsection{PID Evaluation Protocol}

The default run intentionally did not execute fusion. A valid fusion input must
contain a dynamic score, \texttt{row\_id}, session identifier, sequence order, report
path, and SHA-256 digest for each observation. Metadata must declare the signal
as observed, document the score construction and collection policy, and confirm
that the score did not use the target label. Validation and test sessions must be
disjoint.

When such evidence becomes available, the runner joins each dynamic observation
to saved final-round active-learning probabilities. It tunes fixed, P, PI, and
PID rules only on validation, freezes the selected settings, evaluates on test,
and saves all fused predictions. Static-only, Dynamic-only, Average, and untuned
reference PID are included. The runner reports score-replay milliseconds per row
but labels this value as excluding collection, model inference, and I/O. No
fusion metric appears in Section~\ref{sec:results} because the required input is
absent.

\subsection{Independent Replay Verification}

The verification command reloads every active-learning checkpoint and
recomputes probabilities, metrics, and selected queries from saved row indices.
For fidelity, it reloads the model, training baseline, mask indices, and SHAP
inputs, reconstructs each masked matrix, and recomputes predictions and geometry.
The review package reports successful replay of 150 active-learning checkpoints,
120,000 active-learning prediction records, 12,500 query records, 63 mask
batches, and 25,200 masked prediction records. Verification establishes
computational consistency with the saved artifact; it cannot authenticate the
external origin of the CSV.
